\section{Excluding the remaining small standard alternating tensor cores}
\label{sec:alternating-standard-tensors}

This section closes the four necessary tensor cases left by
HD-ALT-02b.  The key estimate is for nonreal eigenvalues of the
actual rational orthogonal tensor operator before its scalar
phase is applied.  It is proved from traces of every power;
no assertion that scalar convolution preserves a nonreal
eigenspace bound is used.

\paragraph{Review status.}
This complete new author derivation and its entire checker have
passed the root reviewer's independent whole-proof reading.
The root reviewer also freshly reran the exact certificate:
$463$ partitions, $20551$ two-factor tensor spectra and $463$
swap spectra passed in $3.787$ seconds.  This rerun is
distinguished from a separate arithmetic implementation.

\subsection*{D-AST-01. Traces of all tensor-wreath powers}

Let $n\ge3$, let $W=\operatorname{Std}_n$ be the deleted
permutation module over $\mathbf Q$, and write $a=n-1$.
For $t\ge1$, let $\mathcal R_{n,t}$ be the image on
$V=W^{\otimes t}$ of the tensor-wreath action of
$S_n\wr S_t$.  It consists of tensor products of standard
permutation matrices followed by permutations of tensor
factors.  Put $e=a^t$.

\begin{lemma}[D-AST-01: lower trace bound for every power]
\label{D-AST-01}
For every $T\in\mathcal R_{n,t}$ and every integer $r$,
\[
                         \operatorname{tr}(T^r)\ge-e/a.
                                                               \tag{AST.1}
\]
Every nonreal eigenvalue of $T$ has multiplicity at most
\[
                              \frac{a+1}{3a}e.             \tag{AST.2}
\]
\end{lemma}

\begin{proof}
For a permutation $\pi\in S_n$ one has
\[
          \operatorname{tr}(\operatorname{Std}_n(\pi))
                       =\#\operatorname{Fix}(\pi)-1
                       \in\{-1,0,1,\ldots,a\}.            \tag{AST.3}
\]
For a tensor operator followed by a permutation of the factors,
contraction of the tensor-basis indices in its trace gives one
factor for each cycle of that factor permutation.  The factor
is the trace of the product of the standard matrices around
the cycle.  That product is another standard permutation
matrix.  Hence, if there are $c$ cycles, its trace is a product
of $c$ integers from (AST.3).

If this product is negative, one of its factors is $-1$;
the absolute values of the remaining $c-1$ factors are at
most $a$.  Thus the product is at least $-a^{c-1}$, which
is at least $-a^{t-1}=-e/a$.  If the product is nonnegative
the same bound is immediate.  Every power $T^r$ belongs to
the same finite tensor-wreath image, so the argument proves
(AST.1) for every integer $r$, not just for $r=1$.

The matrices in $\mathcal R_{n,t}$ are real and orthogonal
for the tensor product of the standard positive forms.
Thus all eigenvalues have absolute value one, and the
multiplicities of an eigenvalue $\zeta$ and its complex
conjugate agree.  Suppose $\zeta$ is nonreal, with order
$d\ge3$ and multiplicity $m$.  Some power satisfies
\[
                         \operatorname{Re}(\zeta^r)\le-1/2.
\]
Indeed, writing $\zeta=\exp(2\pi i b/d)$ with $(b,d)=1$,
choose $r$ so that $br\equiv\lfloor d/2\rfloor\pmod d$.
The resulting argument is in $[2\pi/3,\pi]$.
The two original eigenspaces contribute at most $-m$ to
the real trace of $T^r$, and all other $e-2m$ eigenvalues
contribute at most $e-2m$.  Consequently
\[
                      -e/a\le\operatorname{tr}(T^r)
                                  \le e-3m.
\]
This is (AST.2).  The two original eigenspaces are distinct
even when their powers coincide at the real value $-1$, so
that possibility does not affect the estimate.
\end{proof}

\subsection*{D-AST-02. Extraction of the actual scalar phase}

\begin{lemma}[D-AST-02: the full standard tensor normalizer]
\label{D-AST-02}
Let $n\ge8$.  Let $J$ be a finite linear group normalizing
the irreducibly acting ordinary core $A_n^t$ on
$\operatorname{Std}_n^{\otimes t}$.  Then every $g\in J$
has the form
\[
                     g=\lambda T,\qquad
             T\in\mathcal R_{n,t},\quad \lambda\in\mathbf C^*.
                                                               \tag{AST.4}
\]
For the tensor product standard symmetric form, the multiplier
of $g$ is $\lambda^2$.  If that multiplier is nontrivial,
\[
                 \operatorname{rank}(1-g)
                         \ge\frac{2a-1}{3a}e.              \tag{AST.5}
\]
Every other nonidentity $g$ has moving rank at least $e/a$.
\end{lemma}

\begin{proof}
The simple direct factors of the centerless core are intrinsic,
so $J$ permutes them and their full matrix algebras.  Remove
the corresponding standard permutation of tensor factors.
Each remaining factor automorphism is an element of
$\operatorname{Aut}(A_n)=S_n$, the same full-automorphism
input used in HD-ALT-01a and HD-ALT-02a.  The standard
$S_n$ matrix implements that automorphism on the factor
module.  After removing all these matrices, the residual
operator centralizes the irreducibly acting entire core and
is scalar by Schur's lemma.  This proves (AST.4), including
all outer actions and all factor permutations.  It is an
extraction of the scalar phase from the actual operator
before its nonreal eigenspaces are estimated.

The operator $T$ is rational and preserves the positive
standard tensor form, so the multiplier is $\lambda^2$.
Both $g$ and $T$ have finite order, whence $\lambda$ is a
root of unity.  If $\lambda^2\ne1$, then
$\lambda^{-1}$ is nonreal.  The fixed space of $g$ is the
$\lambda^{-1}$ eigenspace of this particular $T$.
Applying (AST.2) to it gives (AST.5).

For the remaining lower bound, a nonscalar standard factor
matrix has largest eigenspace codimension at least one.
Its fraction passes to a factor-preserving tensor operator
by fixing the eigenvalues in all other factors, giving
$\nu(g)\ge e/a$.  A nontrivial factor-permutation cycle
has the stronger bound $\nu(g)\ge e/4$ by the weighted
constant-word argument, and $a\ge7$ here.  If the actual
matrix is a nonidentity scalar, its moving rank is $e$.
This proves the last assertion in all cases.  Multiplication
by a scalar preserves the largest-eigenspace codimension
$\nu$; it is not used to preserve the nonreal-eigenvalue
bound (AST.2).
\end{proof}

\subsection*{D-AST-03. Closure of the finite alternating endpoint}

\begin{theorem}[D-AST-03: no standard tensor core for $8\le n\le14$]
\label{D-AST-03}
Let an inducing block of an actual rank-nine Weil sixfold
finite-cover source have an irreducibly acting normal ordinary
core $A_n^t$, with representation
$\operatorname{Std}_n^{\otimes t}$, where
$8\le n\le14$ and $t\ge2$.
Suppose its inducing character field contains the imaginary
quadratic original field, as it does for the induced sources
under discussion.  These data cannot occur.

In particular the four remaining cases of HD-ALT-02b are
excluded.  Combined with that theorem, no forced ordinary
alternating core of degree $8\le n\le14$ occurs in its
dominant nonsplit source problem.  The assertion allows all
scalar phases, larger inducing fields, and compatible tensor
permutations.
\end{theorem}

\begin{proof}
Put $a=n-1$ and $e=a^t$.  The unique invariant line of
symmetric forms is normalized by $J$, with multiplier
character $\mu$.  If $\mu$ were trivial, the character of
the finite self-dual representation of $J$ would be real.
That contradicts the inclusion of the original imaginary
quadratic field in its inducing character field.  Thus
$\mu$ is nontrivial.

The actual intermediate-cover Euler identity and period
factorization, already proved for arbitrary inducing fields,
force a rational intermediate base, at least twelve genuine
inertias, and
\[
                 \sum_{j=1}^r\operatorname{rank}(1-g_j)
                           =2e+6,\qquad r\ge12.             \tag{AST.6}
\]
For clarity, genus one would have at most six genuine
branches and a marked-base variation of dimension at most
six, whereas higher genus makes
$6+2e(1-h)$ negative since $e\ge49$.  Thus (AST.6) is a
constraint on the actual period map.

Because the inertias generate $J$ and have product one,
at least two have nontrivial multiplier.  Applying (AST.5)
to two such entries and the last assertion of D-AST-02
to the at least ten others yields
\[
 2e+6\ \ge\ 2\frac{2a-1}{3a}e+10\frac ea
             =\frac{4a+28}{3a}e.
\]
Since $7\le a\le13$, this implies
\[
                          e\le\frac{9a}{14-a}.              \tag{AST.7}
\]
But $t\ge2$ gives $e\ge a^2$, and
\[
              a^2>\frac{9a}{14-a}
             \quad\Longleftrightarrow\quad a(14-a)>9.
\]
On the integer interval $7\le a\le13$ the left-hand
side is at least thirteen.  This is a contradiction.
For the four previously remaining values the exact bounds
are, respectively,
\[
\begin{array}{c|rrrr}
n&11&12&13&14\\\hline
e\text{ upper bound}&45/2&33&54&117\\
e\text{ lower bound}&100&121&144&169.
\end{array}
\]
This proves the theorem.
\end{proof}

\paragraph{Certificate and scope.}
The exact program \texttt{check\_alternating\_standard\_tensor.py}
checks every endpoint inequality.  It also reconstructs the
deleted permutation spectra from all partitions for
$8\le n\le14$, and independently checks (AST.2) on all
two-factor tensor spectra and every two-factor swap spectrum.
The proof for arbitrary $t$ is (AST.1), not an extrapolation
from that finite computation.  Nothing here asserts a bound
for arbitrary shifted convolutions of real spectra or excludes
the ordinary alternating range $n\ge15$.
